\chapter{Requirements and Installation}
\label{ch:requirements}

\section{Supported platform}

Blackship builds and runs on FreeBSD 15.0 or newer. It uses FreeBSD-specific
headers, system calls, ioctls, jail parameters, and network facilities; other
operating systems are not supported build targets.

The minimum host provides:

\begin{itemize}
  \item a FreeBSD 15.0+ kernel with jail support;
  \item the base userland tools, including \command{jail}, \command{jls},
    \command{jexec}, \command{mount}, and \command{ifconfig};
  \item root authority for operations that create or alter host resources; and
  \item network access to a FreeBSD mirror when bootstrapping a release.
\end{itemize}

\subsection{Optional host facilities}

\begin{longtable}{L{0.18\textwidth}L{0.72\textwidth}}
\toprule
Facility & Purpose \\
\midrule
ZFS & Clone provisioning, build-layer cache, snapshots, commit, efficient export,
and migration. \\
PF & Port forwarding, NAT integration, and connection draining. \\
RACCT/RCTL & Memory, CPU, process, file, and custom resource limits; usage sampling. \\
VNET modules & Per-jail network stacks through epair/\command{if\_bridge} or
netgraph. \\
cpuset & Explicit CPU pinning or least-loaded core allocation. \\
\bottomrule
\end{longtable}

RCTL requires the kernel accounting facility at boot. Add the following to
\file{/boot/loader.conf} and reboot before relying on resource policy:

\begin{lstlisting}
kern.racct.enable="1"
\end{lstlisting}

\section{Installing a release binary}

Download the asset appropriate to the host architecture from the project's
release page, extract it, and install the executable on the system path:

\begin{lstlisting}[language=bash]
fetch https://github.com/seuros/blackship/releases/latest/download/blackship-freebsd-amd64.tar.gz
tar xzf blackship-freebsd-amd64.tar.gz
sudo install -m 0555 blackship /usr/local/bin/blackship
blackship --version
\end{lstlisting}

The FreeBSD port under \file{ports/sysutils/blackship} provides the expected
packaging metadata and rc integration for a ports-based deployment.

\section{Installing with Cargo}

Building from source requires a Rust toolchain that satisfies the
\file{Cargo.toml} \command{rust-version}, a C compiler, and the base system
headers. Blackship compiles a small C helper at build time to obtain ioctl and
netgraph constants from the host headers instead of hard-coding ABI values.

\begin{lstlisting}[language=bash]
pkg install rust
cargo install blackship
\end{lstlisting}

For repository development:

\begin{lstlisting}[language=bash]
git clone https://github.com/seuros/blackship.git
cd blackship
cargo build --release
sudo install -m 0555 target/release/blackship /usr/local/bin/blackship
\end{lstlisting}

The default build includes the CLI, telemetry support, and MCP server. Library
consumers can omit them with \command{default-features = false}; Chapter
\ref{ch:automation} covers feature selection.

\section{Checking host capabilities}

The capability report needs neither a manifest nor root:

\begin{lstlisting}[language=bash]
blackship capabilities
blackship --format json capabilities
\end{lstlisting}

It reports the Blackship version, JSON schema version, target OS and
architecture, kernel-derived features, compiled Cargo features, canonical host
tools, and installed privilege escalators. Use it as the first handshake for a
management application.

\section{Data and cache directories}

Defaults follow the XDG base directory conventions:

\begin{longtable}{L{0.28\textwidth}L{0.42\textwidth}L{0.22\textwidth}}
\toprule
Purpose & Default & Override \\
\midrule
Persistent data & \path{~/.local/share/blackship} & \command{XDG\_DATA\_HOME} \\
Bootstrapped releases & \path{<data_dir>/releases} & manifest \\
Download/build cache & \path{~/.cache/blackship} & \command{XDG\_CACHE\_HOME} \\
Manifest & \path{./blackship.toml} & \command{--config} \\
\bottomrule
\end{longtable}

A system-wide installation commonly sets \setting{data\_dir} to
\file{/var/blackship}. Whatever path is chosen must remain consistent across
interactive commands and the supervisor.

\section{Privilege escalation}

Library calls never elevate themselves. The CLI can re-execute privileged
commands through, in order, FreeBSD \command{mdo}, \command{doas}, or
\command{sudo}. The helper path is canonical rather than resolved from an
untrusted \command{PATH}.

\begin{lstlisting}[language=bash]
# Refuse automatic elevation.
blackship --no-escalate up --all

# Elevation may occur, but doas/sudo must not prompt.
blackship --non-interactive up --all
\end{lstlisting}

\command{mdo} authorization is supplied by \command{mac\_do} policy and is
already non-interactive. The CLI resets root identity and XDG environment
variables when it re-executes, preventing privileged state from being written
under the invoking user's home by accident.

\begin{shipwarning}[Bootstrap still needs root]
Even when the destination directory is user-writable, extracting a FreeBSD base
archive needs root to preserve ownership and special permission bits. Do not
work around this by stripping archive metadata; the resulting userland is not a
faithful release.
\end{shipwarning}

\section{Shell completion}

\begin{lstlisting}[language=bash]
blackship completion bash > /usr/local/etc/bash_completion.d/blackship
blackship completion zsh > /usr/local/share/zsh/site-functions/_blackship
blackship completion fish > ~/.config/fish/completions/blackship.fish
\end{lstlisting}
